\documentclass[11pt,a4paper]{article}
\usepackage[a4paper,margin=14mm,top=12mm,bottom=14mm]{geometry}
\usepackage{fontspec}
\setmainfont{Liberation Sans}
\usepackage{polyglossia}
\setdefaultlanguage{french}
\usepackage{amsmath}
\usepackage{xcolor}
\usepackage{graphicx}
\usepackage[most]{tcolorbox}
\usepackage{tabularx,booktabs,array,multirow}
\usepackage{enumitem}
\usepackage{fancyhdr}
\usepackage{lastpage}
\usepackage[hidelinks]{hyperref}
\definecolor{uimm}{HTML}{005677}
\definecolor{cyan}{HTML}{13A2DD}
\definecolor{fond}{HTML}{E6F2FA}
\definecolor{ok}{HTML}{2E7D32}
\definecolor{ko}{HTML}{B3261E}
\definecolor{violet}{HTML}{7B1FA2}
\setlength{\parindent}{0pt}
\setlength{\parskip}{3pt}
\setlist{nosep}
\setlist[itemize]{leftmargin=13pt,label=\textcolor{cyan}{\textbullet}}
\newcommand{\titre}[1]{\vspace{5pt}{\color{uimm}\large\bfseries #1}\par\vspace{-3pt}{\color{cyan}\rule{\linewidth}{1.2pt}}\par\vspace{0pt}}
\tcbset{boite/.style={colback=fond,colframe=uimm,boxrule=0.6pt,arc=2pt,left=6pt,right=6pt,top=4pt,bottom=4pt},
        defi/.style={colback=violet!6,colframe=violet,boxrule=0.8pt,arc=2pt,left=6pt,right=6pt,top=4pt,bottom=4pt,
                     title={\bfseries Défi},coltitle=white,colbacktitle=violet,fonttitle=\small}}
\newcommand{\bandeau}[3]{%
\begin{tcolorbox}[colback=uimm,colframe=uimm,arc=2pt,left=4pt,right=8pt,top=4pt,bottom=4pt]
\begin{minipage}[c]{18mm}\setlength{\fboxsep}{1.2pt}\fcolorbox{white}{white}{\includegraphics[height=16.5mm]{logo_uimm.jpg}}\end{minipage}\hspace{4mm}%
\begin{minipage}[c]{\dimexpr\linewidth-22mm\relax}\color{white}
{\small #1}\\[1pt]{\LARGE\bfseries #2}\\[2pt]{\large #3}\hfill{\small CQPM Technicien qualité}
\end{minipage}
\end{tcolorbox}}
\newcommand{\reponse}[1][1]{\par\vspace{2pt}{\color{gray!60}\foreach \i in {1,...,#1}{\rule{\linewidth}{0.3pt}\par\vspace{9pt}}}}
\usepackage{pgffor}
\pagestyle{fancy}\fancyhf{}
\renewcommand{\headrulewidth}{0pt}
\newcommand{\pied}[1]{\fancyfoot[L]{\footnotesize\color{uimm}#1}\fancyfoot[R]{\footnotesize\color{uimm}Page \thepage{} sur \pageref{LastPage}}}

\usepackage{tikz}
\usepackage{longtable}
\renewcommand{\bandeau}[3]{%
\begin{tcolorbox}[colback=uimm,colframe=uimm,arc=2pt,left=4pt,right=8pt,top=4pt,bottom=4pt]
\begin{minipage}[c]{18mm}\setlength{\fboxsep}{1.2pt}\fcolorbox{white}{white}{\includegraphics[height=16.5mm]{logo_uimm.jpg}}\end{minipage}\hspace{4mm}%
\begin{minipage}[c]{\dimexpr\linewidth-22mm\relax}\color{white}
{\small #1}\\[1pt]{\LARGE\bfseries #2}\\[2pt]{\large #3}
\end{minipage}
\end{tcolorbox}}
% en-tête d'une fiche : code, titre, question, bloc, durée, organisation, niveau
\newcommand{\entete}[7]{%
\begin{tcolorbox}[colback=uimm,colframe=uimm,arc=2pt,left=4pt,right=6pt,top=4pt,bottom=4pt]
\begin{minipage}[c]{18mm}\setlength{\fboxsep}{1.2pt}\fcolorbox{white}{white}{\includegraphics[height=16.5mm]{logo_uimm.jpg}}\end{minipage}\hspace{3mm}%
\begin{minipage}[c]{\dimexpr\linewidth-58mm\relax}\color{white}\raggedright
{\small Base SPC · #4}\\[1pt]{\Large\bfseries #1 · #2}\\[2pt]{\itshape #3}
\end{minipage}\hfill
\begin{minipage}[c]{33mm}\color{white}\small\raggedleft
Durée : \textbf{#5}\\ \textbf{#6}\\ Niveau #7 sur 3
\end{minipage}
\end{tcolorbox}}
\newcommand{\tagE}{{\setlength{\fboxsep}{1.5pt}\colorbox{cyan}{\color{white}\scriptsize\bfseries Ellistat}}}
\newcommand{\tagW}{{\setlength{\fboxsep}{1.5pt}\colorbox{uimm}{\color{white}\scriptsize\bfseries Page web}}}
\tcbset{pred/.style={colback=orange!6,colframe=orange!80!black,boxrule=0.6pt,arc=2pt,left=6pt,right=6pt,top=3pt,bottom=1pt,
                    title={\bfseries Je prédis, avant de commencer},coltitle=white,colbacktitle=orange!80!black,fonttitle=\small},
        retenir/.style={colback=fond,colframe=uimm,boxrule=0.6pt,arc=2pt,left=6pt,right=6pt,top=3pt,bottom=3pt,
                    title={\bfseries À retenir},coltitle=white,colbacktitle=uimm,fonttitle=\small},
        plus/.style={colback=violet!6,colframe=violet,boxrule=0.6pt,arc=2pt,left=6pt,right=6pt,top=3pt,bottom=3pt,
                    title={\bfseries Pour aller plus loin},coltitle=white,colbacktitle=violet,fonttitle=\small},
        act/.style={colback=uimm!10,colframe=uimm,boxrule=0.5pt,arc=2pt,left=5pt,right=5pt,top=2pt,bottom=2pt,before skip=7pt,after skip=3pt},
        verif/.style={colback=orange!6,colframe=orange!80!black,boxrule=0.6pt,arc=2pt,left=6pt,right=6pt,top=3pt,bottom=3pt,
                    title={\bfseries À vérifier sur poste avant la première séance},coltitle=white,colbacktitle=orange!80!black,fonttitle=\small}}
\newcommand{\fichier}[1]{\texttt{\small #1}}
\newcommand{\menu}[1]{\textsf{\textbf{#1}}}
\newcommand{\champ}[1]{\rule{#1}{0.3pt}}

\usepackage{amssymb}
\newcommand{\cb}{$\square$\,}
\newcommand{\etape}[2]{\par\vspace{9pt}{\setlength{\fboxsep}{3pt}\colorbox{uimm}{\color{white}\bfseries\large\ #1\ }}\hspace{4pt}{\color{uimm}\large\bfseries #2}\par\vspace{2pt}}
\newcommand{\motnouveau}[1]{\begin{tcolorbox}[colback=violet!6,colframe=violet,boxrule=0.6pt,arc=2pt,left=6pt,right=6pt,top=3pt,bottom=3pt,title={\bfseries Mot nouveau},coltitle=white,colbacktitle=violet,fonttitle=\small]#1\end{tcolorbox}}
\newcommand{\chiffre}[1]{\begin{tcolorbox}[colback=orange!8,colframe=orange!80!black,boxrule=1pt,arc=3pt,left=8pt,right=8pt,top=5pt,bottom=5pt]#1\end{tcolorbox}}

\pied{Journée du 25/11/2026 · Trame terrain 5 · La cale dans la main}
\begin{document}
\thispagestyle{fancy}
\bandeau{Journée du 25/11/2026 · TP 5 sur 5 · durée 25 min · en binôme}{La cale dans la main}{La chaleur des doigts se voit-elle au micromètre ?}
\vspace{-2pt}
\begin{tcolorbox}[boite,left=5pt,right=5pt,top=2pt,bottom=2pt]\small\textbf{Il vous faut.} un micromètre, la cale la plus longue qu'il peut mesurer (20 mm au 0-25, 30 mm au 25-50), un chronomètre, Ellistat.\end{tcolorbox}

\etape{1}{Je prédis}
Une cale d'acier tenue 3 minutes dans la main fermée s'allonge de (1 µm = 0,001 mm) :\par\smallskip\cb 0,1 µm \quad \cb 3 µm \quad \cb 30 µm \quad \cb pas du tout
\etape{2}{Je mesure}
\begin{enumerate}[label=\textbf{\alph*.},leftmargin=16pt,itemsep=1pt]
\item Mesurer \textbf{3 fois} la cale à la température de la salle. Moyenne = \textbf{référence} : \champ{25mm} mm
\item Garder la cale \textbf{3 minutes} dans la main fermée.
\item La poser et la mesurer \textbf{aussitôt}, puis toutes les 30 secondes pendant 5 minutes, sans la reprendre en main.
\end{enumerate}
\par\vspace{2pt}{\setlength{\tabcolsep}{2pt}\renewcommand{\arraystretch}{1.1}\begin{tabular}{|l|p{11.3mm}|p{11.3mm}|p{11.3mm}|p{11.3mm}|p{11.3mm}|p{11.3mm}|p{11.3mm}|p{11.3mm}|p{11.3mm}|p{11.3mm}|p{11.3mm}|}\hline
\textbf{Temps (s)} & \textbf{0} & \textbf{30} & \textbf{60} & \textbf{90} & \textbf{120} & \textbf{150} & \textbf{180} & \textbf{210} & \textbf{240} & \textbf{270} & \textbf{300} \\ \hline
\rule{0pt}{8mm}\textbf{Mesure} &  &  &  &  &  &  &  &  &  &  &  \\ \hline
\end{tabular}}\par
\etape{3}{Je saisis dans Ellistat}
Dans une grille, une colonne \textbf{Chaleur} : taper les 11 mesures \textbf{dans l'ordre du temps}. \menu{Stats Desc.}, onglet \menu{Détail}, cliquer sur la colonne : le graphique « Moyenne » montre les points dans l'ordre. Dessiner ci-dessous l'allure de la courbe :
\par\vspace{2pt}\fbox{\rule{0pt}{20mm}\rule{0.97\linewidth}{0pt}}
\etape{4}{Je calcule un seul chiffre}
\chiffre{Allongement = première mesure après la main $-$ référence\hfill = \champ{25mm} µm\hfill\ \par\smallskip\small Repère : l'acier s'allonge d'environ 11,5 µm par mètre et par degré (valeur usuelle, à confirmer chez le fabricant des cales). Pour 30 mm, cela fait environ 0,35 µm par degré.}
\etape{5}{Je décide}
Au poste, avant de mesurer une pièce que l'on vient de tenir en main, j'attends \champ{14mm} minutes.\par\smallskip
Une précaution à retenir :\par\vspace{1pt}\reponse[1]
\motnouveau{\textbf{Température de référence} : les cales et les pièces sont définies à 20 °C. Au millième, une pièce plus chaude est une pièce plus longue.}

\end{document}
